\documentclass[10pt,twocolumn]{article}
\usepackage[utf8]{inputenc}
\usepackage{hyperref}
\usepackage{geometry}
\usepackage{enumitem}
\usepackage{listings}
\usepackage{xcolor}
\usepackage{graphicx}
\usepackage{booktabs}
\usepackage{multirow}
\usepackage{float}
\hypersetup{hidelinks}
\geometry{a4paper, top=0.75in, bottom=0.75in, left=0.55in, right=0.55in}

\lstset{
  basicstyle=\ttfamily\scriptsize,
  breaklines=true,
  frame=single,
  backgroundcolor=\color{gray!10},
  columns=fullflexible
}

\title{Porting Deft to Simulated CXL:\\Benchmarking, Stress Testing, and What Broke\\[4pt]
{\small CS6204 Advanced Topics In Systems: Disaggregated Memory \quad
Instructor: Sam H. Noh}}
\author{Yazeed Alharthi}
\date{May 2026}

\begin{document}
\maketitle

\noindent\textbf{Artifact links.}
Source code: \url{https://github.com/yazeed1s/deft}.
Zipped results: \url{https://drive.google.com/file/d/1X4PLVyHkwPTo89u8wIeRD-bmwHYo1Vcb/view?usp=sharing}.

%──────────────────────────────────────────────────────────────
\section{Introduction}
%──────────────────────────────────────────────────────────────

This is the Part~2 report for the CS6204 course project on Deft, a B+tree
index for disaggregated memory~\cite{deft2025eurosys}.
Part~1 was about bringing up the RDMA benchmarking infrastructure on CloudLab
and getting a reproducible end-to-end run---that took four reservation cycles
and most of its words were about infrastructure failures.
Part~2 does three new things: (1)~ports Deft's entire communication layer from
RDMA to a simulated CXL backend using POSIX shared memory, (2)~runs a full
RDMA vs.\ CXL comparison campaign, and (3)~stress-tests the system through a
five-phase breakage suite covering 422 parameter combinations.

The CXL port is the main code contribution.
The breakage suite is the main experimental contribution.
This report covers both, in that order, after establishing the necessary
background and describing the experimental methodology.

%──────────────────────────────────────────────────────────────
\section{Background}
%──────────────────────────────────────────────────────────────

\subsection{Disaggregated Memory and RDMA}

The class focused on systems where compute and memory are physically separated
and connected by high-speed fabrics.
One-sided RDMA lets a compute node read or write remote memory without
involving the remote CPU.
This changes the design space for data structures: lock protocols, consistency
mechanisms, and I/O granularity all need rethinking when every operation is a
network round trip.

\subsection{What Deft Is}

Deft~\cite{deft2025eurosys} is a B+tree for disaggregated memory published at
EuroSys~'25.
Compute nodes run all tree logic; the memory node is a dumb pool that responds
only to \texttt{MALLOC} and \texttt{NEW\_ROOT} control messages.
Its three main contributions address three specific bottlenecks.

\textit{Hash-based leaf nodes.}
Standard RDMA forces reading a full 1~KB node even if only one key is needed.
Deft's leaves use a hash layout so a single 128-byte bucket pair can be read in
one RDMA verb.

\textit{Segmented internal nodes.}
Internal nodes are split into sub-segments; only the relevant segment is
fetched, reducing per-traversal bandwidth.

\textit{SX (Shared-Exclusive) lock.}
A custom lock encoded in a single 64-bit word.
Four 16-bit counters (S-ticket, S-current, X-ticket, X-current) are updated
with a \emph{masked} RDMA FAA that increments individual lanes without carries
spilling across boundaries.
This is the most performance-sensitive primitive in the system, and it is
exactly where the first CXL correctness bug appeared.

\subsection{Why CXL}

CXL.mem provides cache-coherent load/store access to remote memory.
The CPU can dereference a pointer into CXL-attached memory as if it were local
DRAM, with no queue pair, no completion queue, and no memory registration.
A CXL-backed Deft should run the same B+tree logic over either transport;
the cost profile of each operation changes, but no tree code changes.
That is the hypothesis this project tests.

%──────────────────────────────────────────────────────────────
\section{CXL Port: Design and Implementation}
%──────────────────────────────────────────────────────────────

\subsection{Original Code Structure and DSM Layer}

Before the CXL changes, the original Deft repo is organized in a layered way.
\texttt{src/Tree.cpp} and related tree files contain B+tree logic.
\texttt{include/} and \texttt{src/} also contain the communication layer
(\texttt{dsm\_client.*}, \texttt{dsm\_server.*}, RDMA connection classes).
\texttt{test/} has server/client entry binaries, and \texttt{script/} has
automation for setup, campaign runs, and plotting.

In this codebase, DSM means the remote disaggregated shared-memory pool that
the tree accesses through global offsets (global addresses), not raw pointers.
\texttt{DSMServer} owns and exports that memory pool.
\texttt{DSMClient} is the client-side API used by tree code to do reads, writes,
atomics, and allocator RPC calls on that pool.
So the tree layer talks to DSM API, and DSM API talks to RDMA (or now CXL).

\subsection{Key Abstraction: DSMClient}

The Tree layer calls only \texttt{DSMClient} methods.
It never touches \texttt{ibv\_post\_send}, never sees a queue pair.
This means the entire RDMA stack lives below \texttt{DSMClient}'s public
interface---a CXL implementation just provides the same interface with a
different backing.
Zero tree lines change.
So the transition step is: keep the same \texttt{DSMClient} method contracts,
but replace the RDMA transport primitives used inside those methods.
The next subsection explains the concrete replacement used in this project
(POSIX shared memory regions and local atomics).

\subsection{Simulation via POSIX Shared Memory}

We simulate CXL using POSIX shared memory (\texttt{shm\_open} +
\texttt{mmap MAP\_SHARED}).
Both server and client map the same \texttt{/dev/shm} file; loads and stores
are ordinary cache-coherent memory operations visible to both processes.
Three named regions are used:

\begin{itemize}[nosep]
  \item \texttt{/deft\_dsm} --- the B+tree data pool (replaces the
    RDMA-registered huge-page region)
  \item \texttt{/deft\_lock} --- SX lock state (replaces Mellanox on-chip
    device memory)
  \item \texttt{/deft\_rpc} --- message queues for MALLOC/NEW\_ROOT
    (replaces UD QP send/recv)
\end{itemize}

\subsection{Compile-Time Backend Selection}

The transport is selected at build time with no runtime overhead:

\begin{lstlisting}
option(USE_CXL "Simulated CXL backend" OFF)
if(USE_CXL)
  add_definitions(-DUSE_CXL)
  set(LINKS_FLAGS "... -lrt")
else()
  add_definitions(-DUSE_RDMA)
  set(LINKS_FLAGS "... -libverbs")
endif()
\end{lstlisting}

Unused transport sources are excluded from the build entirely.
Table~\ref{tab:verbmap} shows how each RDMA verb maps to a CXL primitive.

\begin{table}[H]
\centering
\footnotesize
\caption{RDMA verb to CXL primitive mapping.}
\label{tab:verbmap}
\begin{tabular}{ll}
\toprule
RDMA operation & CXL implementation \\
\midrule
\texttt{ibv\_post\_send(READ)} + \texttt{poll\_cq} & \texttt{memcpy} + acquire fence \\
\texttt{ibv\_post\_send(WRITE)} + \texttt{poll\_cq} & \texttt{memcpy} + release fence \\
\texttt{ATOMIC\_CMP\_AND\_SWP} & \texttt{compare\_exchange\_strong} \\
\texttt{ATOMIC\_FETCH\_AND\_ADD} & \texttt{fetch\_add} \\
Masked FAA (SX lock) & Software CAS retry loop \\
\texttt{ibv\_poll\_cq} & No-op (ops are synchronous) \\
UD send/recv (RPC) & Shared-memory ring buffer \\
\bottomrule
\end{tabular}
\end{table}

The masked FAA row is the critical one.
ConnectX-5 executes it as a single atomic NIC operation.
Under CXL it becomes a software CAS loop that retries until it can atomically
update only the targeted 16-bit lane.
This is correctness-preserving but has fundamentally different scaling
behavior under high write concurrency, which the results make visible.

\subsection{Implementation Scale and Bugs Found}

The full delta from upstream Deft is 89 files changed, 8195 insertions,
3313 deletions.
The source-only delta (CMake, \texttt{include/}, \texttt{src/},
\texttt{script/}) is 28 files, 2207 insertions, 198 deletions.
Two correctness bugs were found during bring-up, both fixed before any
benchmark data was collected.

\textit{Bug 1: masked FAA semantics.}
\texttt{FaaBound} and \texttt{FaaDmBound} were initially implemented as plain
64-bit \texttt{fetch\_add}, ignoring the \texttt{log\_sz} and \texttt{mask}
parameters.
The SX lock uses \texttt{XS\_LOCK\_FAA\_MASK = 0x8000800080008000} to update
four 16-bit lanes without carry propagation between lanes.
Plain addition corrupts the lock state.
Symptom: \texttt{Deadlock [...] lock upgrade ...} followed by
\texttt{exit(-1)} (code~255).
Fix: CAS-loop emulation that respects lane boundaries defined by the mask.

\textit{Bug 2: multi-client RPC queue indexing.}
With more than one CXL client process, request queues collided because the
layout keyed queues only by (thread, directory), omitting the client index.
The server also undersized the RPC region.
Fix: queues keyed by (client\_id, thread, directory); corrected size formula
published to memcached.

%──────────────────────────────────────────────────────────────
\section{Experimental Setup}
%──────────────────────────────────────────────────────────────

\subsection{Hardware: Paper vs.\ My Setup}

Table~\ref{tab:setup} compares the original paper's evaluation cluster with the
hardware used here.
The main differences are CPU architecture (Intel vs.\ AMD), fabric type
(InfiniBand vs.\ Ethernet), number of compute nodes (9 vs.\ 5), and the
absence of real CXL hardware (replaced by localhost shared memory).

\begin{table}[H]
\centering
\footnotesize
\caption{Paper cluster vs.\ my setup.}
\label{tab:setup}
\begin{tabular}{lll}
\toprule
Item & Paper & My setup \\
\midrule
Node type & 10$\times$ custom server & 6$\times$ CloudLab d6515 \\
CPU & Intel Xeon 6240M & AMD EPYC 7452 (32c) \\
RAM / node & 96~GB & 125~GiB \\
NIC & ConnectX-5, 100G IB & ConnectX-5~Ex, 100G Eth \\
NUMA nodes & 2 & 1 \\
Fabric & InfiniBand & CloudLab Ethernet \\
Max compute nodes & 9~CN & 5~CN \\
Max threads & 1800 & 150 \\
CXL hardware & --- & POSIX shm on localhost \\
\texttt{/dev/shm} size & --- & 63~GB \\
OS / kernel & Ubuntu 18.04 / 4.15 & Ubuntu 18.04 / 4.15 \\
\bottomrule
\end{tabular}
\end{table}

\subsection{Topology Change from Part 1}

Part~1 used Clemson \texttt{r650} nodes (Intel Xeon, ConnectX-6~Dx).
Part~2 moved to Utah \texttt{d6515} nodes (AMD EPYC, ConnectX-5~Ex).
The move was forced by node availability.
The AMD IOMMU on \texttt{d6515} required an additional fix---passthrough mode
(\texttt{iommu=pt amd\_iommu=pt} in GRUB)---before 16~GB RDMA memory
registration would succeed.
This change is permanent across reboots and is discussed further in
Section~\ref{sec:debugging}.

\subsection{Fairness Caveats}

The RDMA vs.\ CXL comparison in this project is a \emph{transport-path
comparison}, not a hardware-equivalent benchmark.
Several important asymmetries exist.

RDMA runs cross a real 100~Gb network between separate machines; CXL runs are
\texttt{memcpy} and atomics on \texttt{/dev/shm} with no network at all.
Latency differences reflect real protocol cost differences plus the wire, not
protocol cost alone.

Masked FAA is a single NIC hardware atomic under RDMA; it is a software CAS
retry loop under CXL.
This asymmetry makes the two paths fundamentally different under high write
concurrency, and it is the main reason the results diverge at scale.

All five CXL client processes run on the same physical machine as the server,
sharing CPU, LLC, and memory bus.
Real CXL deployments would use separate compute nodes connected via a CXL
switch.

%──────────────────────────────────────────────────────────────
\section{Methodology and Test Corpus}
\label{sec:methodology}
%──────────────────────────────────────────────────────────────

\subsection{How a Single Run Works}

Every benchmark job is driven by \texttt{run\_bench.py}.
It reads topology and binary paths from \texttt{global\_config.yaml}, launches
server and client processes via paramiko SSH (or localhost for CXL), waits for
completion, and parses throughput and latency from the client log.
Resource usage (CPU\%, RSS) is sampled every 2~seconds and averaged.
One run produces one row in a CSV file containing: mode, read\_ratio, zipf,
total\_threads, key\_space, repeat, final\_tp\_mops, final\_lat\_us,
cluster\_cpu\_avg\_pct, cluster\_rss\_avg\_mb, status, and exit\_code.

\texttt{run\_campaign.py} wraps \texttt{run\_bench.py} in a Cartesian-product
loop over (modes, read\_ratios, zipfs, repeats).
\texttt{run\_comparison.sh} calls \texttt{run\_campaign.py} twice---once for
RDMA, once for CXL---with the same grid, then merges results and calls the
plotting scripts.
\texttt{run\_breakage\_plan.sh} calls \texttt{run\_campaign.py} once per phase
per transport.

\subsection{Result Sets}

Three result sets were produced at different times.
The baseline comparison set (160 rows, 0 failures) was produced by a single
invocation of \texttt{run\_comparison.sh --ycsba-full --force-hugepage} with
\texttt{CXL\_CLIENT\_COUNT=5}.
The \texttt{--ycsba-full} flag sets the grid to
modes $\in$ \{smoke, small, mid, big\},
zipf $\in$ \{0.0, 0.5, 0.8, 0.99\},
rr = 50, repeats = 5:
$4 \times 4 \times 1 \times 5 = 80$ jobs per transport.
This is the dataset used for all baseline figures (B1--B6).

The stress phases A, B, and C produced 270 rows across multiple CSV files.
Phases D and E produced 304 more rows.
Combined, the full corpus is 734 rows with zero failures.

\subsection{Data Provenance and Corpus Mapping}

To make the corpus auditable, Table~\ref{tab:corpusmap} maps each result folder
to the figures and sections that use it.
The split is intentional: baseline and stress were run in different batches with
different goals.

\begin{table}[H]
\centering
\footnotesize
\caption{Result-folder provenance and where each subset is used.}
\label{tab:corpusmap}
\begin{tabular}{lrrp{3.2cm}}
\toprule
Folder & Rows & Failures & Used by \\
\midrule
\texttt{deft-resultsv4} & 160 & 0 & Baseline B1--B6 (Section~\ref{sec:baseline}) \\
\texttt{deft-resultsv6} & 270 & 0 & Stress A/B/C figures (S1--S4) \\
\texttt{deft-resultsv7} & 304 & 0 & Stress D/E figures (S5--S10) \\
\midrule
\textbf{Total} & \textbf{734} & \textbf{0} & X1 bridge uses baseline (v4) + stress (v6+v7) \\
\bottomrule
\end{tabular}
\end{table}

\subsection{Five Stress Phases Explained}

Table~\ref{tab:phases} gives the concise view.

\begin{table*}[t]
\centering
\footnotesize
\caption{Breakage suite: phase definitions, fixed and varying parameters, run counts.}
\label{tab:phases}
\begin{tabular}{llp{3.8cm}p{3.8cm}rr}
\toprule
Phase & Goal & Fixed parameters & Varying parameters & Per transport & Total \\
\midrule
A & Baseline across modes & rr=50 & mode $\in$ \{smoke,small,mid\},
  zipf $\in$ \{0.0, 0.8, 0.99\}, 3 repeats & 27 & 54 \\
B & rr/zipf stress & mode $\in$ \{small,mid\} & rr $\in$ \{10,50,90\},
  zipf $\in$ \{0.0,0.8,0.99,0.999\}, 2 repeats & 48 & 96 \\
C & Thread/client scaling & key\_space=10M & tpc $\in$ \{1,4,8,16,30\},
  3 rr/zipf pairs, 2 repeats; CXL also sweeps clients $\in$ \{1,3,5\}
  & 30~RDMA / 90~CXL & 120 \\
D & Keyspace isolation & rr=50, zipf=0.99, tpc=8 & keyspace $\in$
  \{1K,100K,1M,10M,100M,400M\}, 2 repeats & 12 & 24 \\
E & Full interaction surface & tpc=8, 1 repeat & rr $\in$ \{0,10,50,90\},
  zipf $\in$ \{0.0,0.8,0.99,0.999\}, keyspace $\in$ \{1K,1M,100M,400M\}
  & 64 & 128 \\
\midrule
\multicolumn{4}{r}{Grand total (phases A--E)} & 181 RDMA / 241 CXL & 422 \\
\bottomrule
\end{tabular}
\end{table*}

\textbf{Phase A} establishes a baseline with multiple modes and Zipf values at
rr=50.
Three repeats give IQR error bars.
It serves mainly as a sanity check that numbers are stable.

\textbf{Phase B} keeps mode fixed at small and mid and sweeps the full
rr $\times$ zipf grid.
The goal is to understand how read/write ratio and skew interact.
rr=0 is write-only; rr=90 is almost read-only.
zipf=0.999 is extreme hot-key skew.

\textbf{Phase C} is the thread-scaling phase.
It fixes key\_space at 10M and sweeps threads-per-client from 1 to 30.
For RDMA, total threads = tpc $\times$ 5~CN $\in$ \{5, 20, 40, 80, 150\}.
For CXL, client count is also swept ($\in$ \{1, 3, 5\}), so CXL covers
total thread counts $\in$ \{1, 3, 4, 5, 8, 12, 16, 20, 24, 30, 40, 45, 48,
80, 90, 150\}.
The CXL client-count dimension isolates whether performance depends on
total threads, threads-per-client, or how clients are distributed.

\textbf{Phase D} isolates keyspace effects by holding everything else
constant (rr=50, zipf=0.99, tpc=8) and varying only key\_space across
six orders of magnitude.

\textbf{Phase E} is the most comprehensive phase.
It sweeps the three-way interaction of rr, zipf, and keyspace simultaneously
with tpc=8 fixed and one repeat per point.
This answers: is the read-ratio sensitivity from phase B modulated by keyspace
or Zipf, and are there interaction effects the earlier single-axis sweeps missed?

%──────────────────────────────────────────────────────────────
\section{Artifact Reproducibility Runbook}
\label{sec:runbook}
%──────────────────────────────────────────────────────────────

This section gives the minimal commands to reproduce the exact datasets used in
this report.
All commands are run from \texttt{/deft\_code/deft/script} on \texttt{mn0}.

\subsection{Baseline Comparison Corpus (160 rows)}

\begin{lstlisting}
cd /deft_code/deft/script
PYTHONUNBUFFERED=1 CXL_CLIENT_COUNT=5 RNIC_ID=0 \
  ./run_comparison.sh --ycsba-full --force-hugepage
\end{lstlisting}

This executes 80 RDMA jobs + 80 CXL jobs and auto-generates merged comparison
plots.
Output roots are:
\texttt{result/rdma-campaign-*},
\texttt{result/cxl-campaign-*},
\texttt{result/comparison-*}.

\subsection{Breakage Corpus (422 rows)}

\begin{lstlisting}
cd /deft_code/deft/script
PYTHONUNBUFFERED=1 CXL_CLIENT_COUNT=5 RNIC_ID=0 \
  ./run_breakage_plan.sh --phase-a --phase-b --phase-c \
    --phase-d --phase-e
\end{lstlisting}

This executes all five stress phases (A--E) and auto-generates per-phase and
aggregate plots.
Output roots are \texttt{result/rdma-phase*}, \texttt{result/cxl-phase*}, and
the consolidated stress plots from \texttt{plot\_breakage\_report.py}.

%──────────────────────────────────────────────────────────────
\section{Baseline Results}
\label{sec:baseline}
%──────────────────────────────────────────────────────────────

\subsection{Throughput by Mode and Zipf}

\begin{figure}[H]
\centering
\includegraphics[width=\columnwidth]{images/final/B1_tp_mode_zipf.png}
\caption{B1: Throughput (Mops/s) by (mode, Zipf), rr=50. CXL leads at every
point; advantage shrinks from 13$\times$ at smoke to 1.2$\times$ at big as
software-FAA ceiling is approached. Both transports nearly Zipf-insensitive.}
\label{fig:b1}
\end{figure}

Figure~\ref{fig:b1} groups throughput by (mode, zipf) across all 160 baseline
runs.
Table~\ref{tab:baseline} shows the exact median values.
CXL leads at every mode, but the margin shrinks dramatically: 13$\times$ at
smoke (2 total threads) to 1.2$\times$ at big (150 threads).
This convergence is the first clear signal of CXL saturation.

\begin{table}[H]
\centering
\footnotesize
\caption{RDMA vs.\ CXL baseline medians (rr=50, zipf=0.99, 5~repeats).}
\label{tab:baseline}
\begin{tabular}{lrrrr}
\toprule
Mode & RDMA TP & RDMA Lat & CXL TP & CXL Lat \\
 & (Mops/s) & ($\mu$s) & (Mops/s) & ($\mu$s) \\
\midrule
smoke & 0.784 & 6.374 & 10.106 & 0.510 \\
small & 3.739 & 6.675 & 17.535 & 1.405 \\
mid & 10.728 & 6.995 & 20.587 & 3.633 \\
big & 17.672 & 8.494 & 20.535 & 7.270 \\
\bottomrule
\end{tabular}
\end{table}

The Zipf axis is nearly flat at rr=50.
At this mix, lock-path contention dominates more than key-distribution effects.

\subsection{Latency by Mode and Zipf}

\begin{figure}[H]
\centering
\includegraphics[width=\columnwidth]{images/final/B2_lat_mode_zipf.png}
\caption{B2: Latency (\textmu{}s) by (mode, Zipf), rr=50. CXL is 12$\times$
lower at smoke and 1.2$\times$ lower at big. RDMA latency is stable at
6--8~\textmu{}s, bounded by the NIC round-trip floor.}
\label{fig:b2}
\end{figure}

Figure~\ref{fig:b2} shows latency.
At smoke, CXL latency is 0.5~\textmu{}s vs.\ RDMA 6.4~\textmu{}s.
At big (150 threads), CXL latency has risen to 7.3~\textmu{}s while RDMA is
8.5~\textmu{}s---the gap has closed to 1.2$\times$.
This is consistent with the software FAA retry cost increasing with thread
count.

\subsection{CPU and RSS}

\begin{figure}[H]
\centering
\includegraphics[width=\columnwidth]{images/final/B3_cpu_mode_zipf.png}
\caption{B3: Cluster CPU (\%) by (mode, Zipf). CXL CPU is 2.5--3$\times$
higher than RDMA at every mode. Flat Zipf axis confirms write-path dominance.}
\label{fig:b3}
\end{figure}

\begin{figure}[H]
\centering
\includegraphics[width=\columnwidth]{images/final/B4_rss_mode_zipf.png}
\caption{B4: Cluster RSS (MB) by (mode, Zipf). CXL RSS is 8--16$\times$ higher
than RDMA because \texttt{/deft\_dsm} is mapped into every client's address
space---a simulation artifact, not a real CXL property.}
\label{fig:b4}
\end{figure}

Figures~\ref{fig:b3} and~\ref{fig:b4} show resource cost.
CXL CPU is $\approx$27,000\% at big mode vs.\ $\approx$11,000\% for RDMA.
Main reasons are busy-polling in the CXL RPC path and software FAA retries;
RDMA offloads both through NIC hardware.

CXL RSS reaches 10,000--24,000~MB vs.\ 600--3,000~MB for RDMA.
The DSM pool is \texttt{mmap}'d into every client process and counts against
each one's RSS.
This is a simulation artifact, not a fundamental CXL property.

One efficiency number worth noting: at smoke, TP/CPU is
$0.784/433 = 0.0018$~Mops/s/CPU\% for RDMA vs.\
$10.1/1056 = 0.0096$ for CXL---CXL is 5.3$\times$ more CPU-efficient despite
the higher absolute cost.
CXL avoids verbs stack overhead, NIC interrupt coalescing, and QP management,
so each CPU cycle does more useful work.

\subsection{Thread Scaling}

\begin{figure}[H]
\centering
\includegraphics[width=\columnwidth]{images/final/B5_tp_threads.png}
\caption{B5: Throughput vs.\ total threads. RDMA scales linearly (0.8 to
17.7~Mops/s). CXL starts 12$\times$ higher but saturates near 20--21~Mops/s
around 75 threads as software-FAA CAS retries exhaust cache-line bandwidth.}
\label{fig:b5}
\end{figure}

\begin{figure}[H]
\centering
\includegraphics[width=\columnwidth]{images/final/B6_lat_threads.png}
\caption{B6: Latency vs.\ total threads. RDMA grows slowly (6.4 to
8.5~\textmu{}s). CXL grows steeply (0.5 to 7.3~\textmu{}s) as each added
thread increases per-operation CAS retry probability.}
\label{fig:b6}
\end{figure}

Figures~\ref{fig:b5} and~\ref{fig:b6} collapse the mode axis into a thread
count axis.
RDMA throughput is nearly linear from 5 to 150 threads.
CXL throughput rises steeply at first then flattens at $\approx$20~Mops/s
around 75 threads.
This is the practical ceiling of software FAA on shared lock words.

The paper's Figure~8 shows RDMA scaling linearly through 1800 threads.
Our RDMA curve shows the same linear shape through 150 threads.
Absolute throughput is lower (17.7 vs.\ $\approx$42~Mops/s), which is expected
given 5~CNs on Ethernet vs.\ 9~CNs on InfiniBand.

%──────────────────────────────────────────────────────────────
\section{Reproduction vs.\ Original Paper}
\label{sec:repro}
%──────────────────────────────────────────────────────────────

This section compares what was reproduced from Deft's EuroSys'25 trends and
what was intentionally different.
Table~\ref{tab:papercompare} is a compact side-by-side summary.

\begin{table*}[t]
\centering
\footnotesize
\caption{Paper vs.\ this report: explicit trend/value comparison and gap reasons.}
\label{tab:papercompare}
\begin{tabular}{p{2.4cm}p{3.5cm}p{3.5cm}p{5.2cm}}
\toprule
Dimension & Paper (Deft EuroSys'25) & This report & Why different / interpretation \\
\midrule
Cluster + fabric & 10 nodes total (1 MN + up to 9 CN), 100G InfiniBand, Intel Xeon 6240M; scales to 1800 threads & 6 nodes total (1 MN + 5 CN), 100G Ethernet, AMD EPYC 7452; scales to 150 threads & Hardware and scale are smaller in this project; absolute values should be lower, but shape/trend is still comparable \\
RDMA scaling trend & Figure~8: throughput grows approximately linearly to 1800 threads; peak around 42~Mops/s & Figure~B5: throughput grows approximately linearly to 150 threads; peak 17.7~Mops/s & Trend is reproduced (linear growth), high-thread regime is not reached because thread ceiling is 150 \\
Latency floor behavior & RDMA path is microsecond-level and bounded by network round-trip and verbs overhead & RDMA remains about 6--8~\textmu{}s; simulated CXL starts at 0.5~\textmu{}s in smoke & Matches expected transport cost split: networked RDMA floor vs.\ local load/store path \\
Workload scope & YCSB + Twitter workloads; includes cross-index comparison (Sherman, dLSM, SMART), reported 2.4--9.5$\times$ gains & YCSB-style campaign + five-phase breakage suite; no cross-index rerun in this project & Project focus is transport porting + robustness, not full competitor re-evaluation \\
Extreme skew corner case & No reported \texttt{zipf=0.999, rr=0} stress point in paper plots & Added in Phase~E: RDMA latency spikes to 29~\textmu{}s, CXL around 3~\textmu{}s & This is an extension beyond reproduction; it reveals a pathological write-skew behavior not visible in the original matrix \\
\bottomrule
\end{tabular}
\end{table*}

In short: the main paper trend that can be fairly compared here is RDMA
thread-scaling shape, and that trend reproduces.
Absolute throughput differs for expected infrastructure reasons.
The stress suite adds new parameter corners that were not in the paper and
exposes additional behavior.

%──────────────────────────────────────────────────────────────
\section{Stress Results}
\label{sec:stress}
%──────────────────────────────────────────────────────────────

All 422 stress runs completed with \texttt{status=ok} and exit code 0.
The \texttt{exit(-1)} deadlock from the masked FAA bug was the only failure
mode ever observed; it was fixed before any stress data was collected.

\subsection{Phase A/B/C: Overall Distributions}

\begin{figure}[H]
\centering
\includegraphics[width=\columnwidth]{images/final/S1_abc_tp_hist.png}
\caption{S1: Throughput distribution across Phase A/B/C runs (270 rows).
RDMA concentrates near 0--3~Mops/s; CXL spreads broadly from 5 to 21~Mops/s.}
\label{fig:s1}
\end{figure}

\begin{figure}[H]
\centering
\includegraphics[width=\columnwidth]{images/final/S2_abc_lat_hist.png}
\caption{S2: Latency distribution across Phase A/B/C runs. RDMA clusters at
6--8~\textmu{}s regardless of parameters. CXL is bimodal: 0.5--1.5~\textmu{}s
at low concurrency and 6--8~\textmu{}s at high concurrency.}
\label{fig:s2}
\end{figure}

Figures~\ref{fig:s1} and~\ref{fig:s2} show throughput and latency histograms
across all Phase~A/B/C runs.
RDMA stays concentrated at lower throughput and 6--8~\textmu{}s latency.
CXL has a wider throughput spread and a bimodal latency shape
(low-concurrency fast mode, high-concurrency contention mode).
This already suggests a crossover thread region; Phase~C quantifies it.

\subsection{Phase C: Thread and Client Scaling}

\begin{figure}[H]
\centering
\includegraphics[width=\columnwidth]{images/final/S3_phasec_tp_threads.png}
\caption{S3: Phase C throughput vs.\ total threads (key\_space=10M). CXL peaks
near 14~Mops/s at 40--50 threads then declines to 11~Mops/s at 150. RDMA grows
monotonically to 10~Mops/s. Curves cross near 100 threads.}
\label{fig:s3}
\end{figure}

\begin{figure}[H]
\centering
\includegraphics[width=\columnwidth]{images/final/S4_phasec_lat_threads.png}
\caption{S4: Phase C latency vs.\ total threads. RDMA stays near 8--9~\textmu{}s,
rising to 15~\textmu{}s at 150. CXL starts near zero and rises steeply to
13~\textmu{}s, crossing RDMA between 80 and 100 threads.}
\label{fig:s4}
\end{figure}

Phase~C is the most revealing phase in the suite.
Figure~\ref{fig:s3} shows CXL throughput exhibiting a peak-and-decline shape:
it rises from $\approx$5.5~Mops/s at 5 threads to $\approx$14~Mops/s at
40--50 threads, then \emph{declines} back to $\approx$11~Mops/s at 150 threads.
RDMA grows monotonically from 0.6 to 10~Mops/s.
This is consistent with contention amplification in software atomics at higher
thread counts.

Figure~\ref{fig:s4} shows the latency crossover quantitatively.
At 1 thread, CXL latency is near 0~\textmu{}s; RDMA is 8.5~\textmu{}s.
By 80--90 threads, they equalize.
At 150 threads, CXL (13.4~\textmu{}s) is still slightly lower than RDMA
(15.3~\textmu{}s).
Practical point: software FAA emulation loses most latency advantage beyond
about 80 concurrent threads.

\subsection{Phase D: Keyspace Isolation}

\begin{figure}[H]
\centering
\includegraphics[width=\columnwidth]{images/final/S5_phased_tp_keyspace.png}
\caption{S5: Throughput vs.\ keyspace (log scale), tpc=8, rr=50, zipf=0.99.
RDMA is stable then drops at 400M keys. CXL is flat then slightly increases at
400M, possibly from hardware prefetching on deeper traversal paths.}
\label{fig:s5}
\end{figure}

\begin{figure}[H]
\centering
\includegraphics[width=\columnwidth]{images/final/S6_phased_cpu_keyspace.png}
\caption{S6: CPU vs.\ keyspace. CXL is nearly flat through 10M then rises to
17,000\% at 400M. RDMA is flat at 2,500--3,400\%.}
\label{fig:s6}
\end{figure}

\begin{figure}[H]
\centering
\includegraphics[width=\columnwidth]{images/final/S7_phased_rss_keyspace.png}
\caption{S7: RSS vs.\ keyspace (log scale). RDMA grows slowly tracking the
local index cache. CXL is non-monotonic: dips at mid-range keyspaces where
Zipf concentrates access, then spikes to 27,000~MB at 400M.}
\label{fig:s7}
\end{figure}

Phase~D holds rr=50, zipf=0.99, and tpc=8 fixed and sweeps keyspace alone.

For RDMA (Figure~\ref{fig:s5}), throughput is flat at 6.3~Mops/s from 1K to
100M keys, then drops to 5.3~Mops/s at 400M.
The drop comes from index cache pressure: the B+tree is deeper at 400M keys,
the client-side cache of internal nodes starts evicting entries, and more RDMA
reads per traversal are needed.

For CXL, throughput is flat at 22.5--23~Mops/s from 1K to 100M, then slightly
\emph{increases} to 24.5~Mops/s at 400M.
The flat region is expected because the simulated CXL pool is in local RAM.
The slight increase at 400M is observed but not fully explained.

Figure~\ref{fig:s7} shows the non-monotonic CXL RSS curve.
At mid-range keyspaces with zipf=0.99, most accesses concentrate on a small
hot set so fewer DSM pages ever fault into physical RAM and RSS is lower.
At 400M keys, even zipf=0.99 touches enough distinct keys to page in a much
larger fraction of the pool.

\subsection{Phase E: Full Interaction Surface}

\begin{figure}[H]
\centering
\includegraphics[width=\columnwidth]{images/final/S8_phasee_tp_rr_zipf.png}
\caption{S8: Throughput vs.\ read ratio, faceted by Zipf (tpc=8). CXL reaches
52--54~Mops/s at rr=90 (6.5$\times$ over RDMA). At rr=0 the advantage is
2.5$\times$. Both transports are Zipf-insensitive at fixed read ratio.}
\label{fig:s8}
\end{figure}

\begin{figure}[H]
\centering
\includegraphics[width=\columnwidth]{images/final/S9_phasee_lat_rr_zipf.png}
\caption{S9: Latency vs.\ read ratio, faceted by Zipf. At zipf=0.999, rr=0,
RDMA latency spikes to 29~\textmu{}s due to extreme hot-key write contention;
CXL stays near 3~\textmu{}s because retries cost nanoseconds not microseconds.}
\label{fig:s9}
\end{figure}

\begin{figure}[H]
\centering
\includegraphics[width=\columnwidth]{images/final/S10_phasee_tp_keyspace_rr.png}
\caption{S10: Throughput vs.\ keyspace, faceted by read ratio. Both transports
are essentially flat. Read ratio dominates; keyspace is secondary at tpc=8.}
\label{fig:s10}
\end{figure}

Phase~E is the most comprehensive: it sweeps read ratio, Zipf, and keyspace
simultaneously with tpc=8 fixed.

Figure~\ref{fig:s8} is the most striking result in the entire corpus.
CXL throughput at rr=90 reaches 52--54~Mops/s---6.5$\times$ over RDMA at the
same read ratio, and more than double CXL's own rr=50 level.
At rr=0 (write-only), the CXL advantage shrinks to 2.5$\times$.
The reason is structural: read-heavy CXL bypasses most retry-loop cost, while
RDMA still pays network round-trip per operation.

Figure~\ref{fig:s9} shows the most extreme behavior in the dataset.
At zipf=0.999, rr=0, RDMA latency spikes to 29~\textmu{}s---nearly three times
higher than at lower Zipf values.
This is hot-lock contention under extreme write skew.
Each RDMA retry pays network RTT, so the tail grows quickly.
CXL stays near 3~\textmu{}s in this case because retries are local-memory
operations, not network trips.
This specific case does not appear in the original paper's evaluation, which
does not test zipf=0.999 at rr=0.

Figure~\ref{fig:s10} confirms that keyspace is inert when tpc=8 is fixed.
Both curves are flat for all read ratios.
The keyspace $\times$ access pattern interaction suggested by Phase~D only
materializes at higher thread counts where a larger working set meaningfully
increases FAA contention.

\subsection{Baseline vs.\ Stress Envelope}

\begin{figure}[H]
\centering
\includegraphics[width=\columnwidth]{images/final/X1_baseline_vs_stress_envelope.png}
\caption{X1: Baseline vs.\ stress medians with stress IQR. Throughput: 14~Mops/s
baseline, 9.5~Mops/s stress median, IQR 5--13~Mops/s. Latency: 6.5~\textmu{}s
baseline, 3.6~\textmu{}s stress median, IQR 1.8--6.7~\textmu{}s.}
\label{fig:x1}
\end{figure}

Figure~\ref{fig:x1} connects the two datasets.
The throughput drop from baseline (14~Mops/s) to stress median (9.5~Mops/s)
reflects the wider parameter space in the stress suite, which includes
write-heavy runs and high-thread configurations where CXL saturates.
Baseline median (6.5~\textmu{}s) is RDMA-dominated.
Stress median (3.6~\textmu{}s) is pulled down by many low-concurrency CXL
runs with sub-microsecond latency.
The IQR of 1.8--6.7~\textmu{}s directly reflects the bimodal CXL distribution
from Figure~\ref{fig:s2}.
The stress envelope is consistent with the explored parameter mix.

%──────────────────────────────────────────────────────────────
\section{Debugging Journey}
\label{sec:debugging}
%──────────────────────────────────────────────────────────────

Table~\ref{tab:bugs} summarizes every significant failure encountered across
both parts of the project.
The subsections below describe the four RDMA infrastructure failures that
required the most work in Part~2.

\begin{table}[H]
\centering
\footnotesize
\caption{Failure log across both project parts.}
\label{tab:bugs}
\begin{tabular}{p{1.8cm}p{2.2cm}p{2.5cm}}
\toprule
Symptom & Root cause & Fix \\
\midrule
QP create errno~22 & Userspace/kernel OFED mismatch & Install \texttt{mlnx-ofed-kernel-dkms} \\
SSH to 10.10.1.x timeout & VLAN destroyed after OFED reset & Rewrite \texttt{net\_heal.sh} to parse ifmap \\
Hugepage preflight fails & Stranded root-owned processes holding pages & sudo-aware killall; no-sudo rule \\
MR ENOMEM at 16~GB & SSH ulimit drops to 65536~KB & \texttt{ulimit -l unlimited} in bash wrapper \\
MR ENOMEM persists & AMD IOMMU SG limit & \texttt{iommu=pt} in GRUB + reboot \\
CXL deadlock (code~255) & Plain 64-bit FAA ignores mask & CAS-loop masked FAA emulation \\
CXL multi-client hang & Queue indexing missing client ID & (client, thread, dir) key + resize RPC region \\
\bottomrule
\end{tabular}
\end{table}

\subsection{RDMA Driver ABI Mismatch}

After moving to Utah \texttt{d6515} nodes, QP creation failed with errno~22.
The device was visible, ports were UP, but both \texttt{ibv\_rc\_pingpong} and
Deft failed the same way.
Cause: \texttt{cloudlab\_setup.sh} installed MLNX\_OFED 4.9 userspace libraries
but not \texttt{mlnx-ofed-kernel-dkms}, so the kernel was running the upstream
Ubuntu \texttt{mlx5\_ib.ko}.
The proprietary experimental verbs passed undocumented struct layouts to the
kernel module and got rejected with \texttt{EINVAL}.
Fix: add \texttt{mlnx-ofed-kernel-dkms} and \texttt{mlnx-ofed-kernel-utils}
to the APT package list and rebuild the kernel module.

\subsection{VLAN Network Disconnection After Driver Reset}

Every \texttt{openibd} restart or reboot destroyed the \texttt{vlan306}
interface and the \texttt{10.10.1.x} experiment addresses.
The original \texttt{net\_heal.sh} reassigned the IP to the bare untagged
physical interface; CloudLab's top-of-rack switches filter untagged traffic.
The rewrite parses \texttt{/var/emulab/boot/ifmap} to extract the correct VLAN
ID, MAC, and IP, then dynamically recreates the tagged subinterface.

\subsection{SSH ulimit Inheritance}

The server's 16~GB MR registration failed with \texttt{ENOMEM} despite 50~GB
of free hugepages.
Paramiko launches a non-interactive SSH session; \texttt{sudo} inside a
non-interactive session resets the locked-memory limit to 65536~KB, which is
about 250,000$\times$ smaller than what we need.
Fix: wrap the launch command as
\texttt{sudo bash -c "ulimit -l unlimited \&\& numactl ... ./server"}.

\subsection{AMD IOMMU Scatter-Gather Limit}

Even with \texttt{ulimit -l unlimited} verified, 16~GB registration still
failed. \texttt{dmesg} showed:

\begin{lstlisting}
0000:41:00.1: IOMMU mapping error in
  map_sg (io-pages: 4194304)
ib_umem_get: failed to map scatterlist
\end{lstlisting}

The AMD IOMMU was rejecting the $\approx$4.19M scatter-gather page mapping
request because IOVA space was exhausted under default protection mode.
Reducing to 8~GB also failed.
At first I did not understand why reducing the registration size didn't help---it
turns out the IOVA limit is not about the total size but about the number of
discrete scatter-gather entries, and both sizes were still exceeding the limit.
Fix: append \texttt{iommu=pt amd\_iommu=pt} to
\texttt{GRUB\_CMDLINE\_LINUX\_DEFAULT}, run \texttt{update-grub}, reboot all
six nodes.
This is a permanent change that survives reboots.

%──────────────────────────────────────────────────────────────
\section{Threats to Validity}
%──────────────────────────────────────────────────────────────

\textbf{Network vs.\ localhost.}
RDMA results include real network latency; CXL results do not.
Latency comparisons read as ``RDMA path including network'' vs.\ ``CXL path on
a single machine.'' The numbers are not hardware-equivalent.

\textbf{Software vs.\ hardware FAA.}
The CAS retry loop is correctness-preserving but has fundamentally different
scaling behavior from NIC hardware atomics.
This is the main reason the two paths diverge at high write concurrency, as
quantified in Figures~\ref{fig:s3}, \ref{fig:s4}, and~\ref{fig:s9}.

\textbf{Single-machine CXL clients.}
Five clients on one machine share CPU, LLC, and memory bus.
Real CXL deployments via a CXL switch would not have this inter-client
contention.

\textbf{Single-server memory pool.}
\texttt{num\_server=1} is assumed throughout.
Multi-server CXL would require a CXL switch and introduces cross-NUMA effects
not modeled here.

\textbf{Thread count range.}
The paper evaluates up to 1800 threads; the maximum here is 150.
The linear RDMA scaling regime is reproduced; the paper's high-concurrency
saturation regime is not.

%──────────────────────────────────────────────────────────────
\section{Lessons Learned}
%──────────────────────────────────────────────────────────────

What I misunderstood at first:
I first treated many failures as RDMA bugs, but several were environment bugs
(driver ABI, VLAN, ulimit, IOMMU).
Uncertainty is not a good feeling.

I thought I was benchmarking Deft; sometimes it felt like Deft was benchmarking
my patience.
This was my first time benchmarking this aggressively, and it was much more
time-consuming than I expected.
Just running the benchmark corpus took around 18 hours end-to-end, not counting
debug/retry cycles.

The hardest part was not writing CXL code; it was making the environment stable
enough to trust measurements.
The main technical win for me is that tree logic stayed unchanged while
transport changed completely.

%──────────────────────────────────────────────────────────────
\section{Conclusion}
%──────────────────────────────────────────────────────────────

This report delivered a complete RDMA-to-CXL transport port for Deft
(20 files modified/created, 46 client methods, zero B+tree logic changes) and
a 734-row evaluation corpus.

This is not a hardware-equivalent RDMA-vs-CXL claim; it is a transport-path
study with simulated CXL.
I reproduced trend shape better than absolute numbers.

Results are consistent across baseline and stress.
At low concurrency, CXL wins strongly by removing RDMA round-trip cost.
At higher concurrency (about $>$80 threads), software FAA retries become the
dominant CXL bottleneck and the advantage narrows.
At read-heavy workloads (rr=90), CXL reaches 52--54~Mops/s; at extreme
write-skew (zipf=0.999, rr=0), RDMA latency spikes to 29~\textmu{}s while CXL
stays near 3~\textmu{}s.

I am still not fully certain about the 400M keyspace CXL throughput bump;
current explanation is a hypothesis, not proof.
By the end, I felt I had more new questions than fully closed answers.

All 422 stress runs completed with no failures.
The full aggressive benchmark execution itself took about 18 hours, which was a
major practical part of this project.

%──────────────────────────────────────────────────────────────
\appendix
\section{Environment Snapshot Command}
%──────────────────────────────────────────────────────────────

For exact environment capture on \texttt{mn0}, this is the command used in the
hardware snapshot note:

\begin{lstlisting}
bash -lc '
echo "===== HOST ====="; hostname -f || hostname; date; uname -a
echo; echo "===== CPU ====="; lscpu
echo; echo "----- NUMA -----"; numactl --hardware 2>/dev/null || true
echo; echo "===== MEMORY ====="; free -h
echo; grep -E "MemTotal|HugePages_Total|HugePages_Free|Hugepagesize" /proc/meminfo
echo; echo "===== STORAGE ====="; lsblk -o NAME,SIZE,TYPE,MOUNTPOINT; df -h
echo; echo "===== NIC / NET ====="; ip -br addr; ip route
echo; lspci | egrep -i "ethernet|network|infiniband|mellanox" || true
echo; echo "===== RDMA ====="; ibv_devices 2>/dev/null || true
echo; ibv_devinfo -l 2>/dev/null || true
'
\end{lstlisting}

%──────────────────────────────────────────────────────────────
\begin{thebibliography}{2}

\bibitem{deft2025eurosys}
J.~Wang, Q.~Wang, Y.~Zhang, and J.~Shu,
``Deft: A Scalable Tree Index for Disaggregated Memory,''
in \emph{Proc.\ EuroSys '25}, Rotterdam, Netherlands, 2025, pp.~886--901.

\bibitem{sherman}
Q.~Wang, Y.~Lu, and J.~Shu,
``Sherman: A Write-Optimized Distributed B$^+$Tree Index on Disaggregated
Memory,''
in \emph{Proc.\ SIGMOD '22}, Philadelphia, PA, USA, 2022, pp.~1033--1048.

\end{thebibliography}

\end{document}
